\section{讨论、结论与研究议程}

最后三张课件把结果拉回研究方法。讲者没有宣称 CoT 已被解释，也没有说 prompt engineering 会永久存在；他列出的是当时能从实验得到的结论、仍需回答的机制问题，以及对学生可参与方向的个人判断。

\subsection{Prompt engineering 会不会消失}

课堂问答区分 easy, well-specified tasks 与 increasingly hard, underspecified tasks。更大模型可能让简单 multiple-choice 对措辞更 robust，但模型越强，人们也会把它用于更复杂目标；此时如何描述约束、提供 evidence、指定 decomposition 与 verification，仍然是 interface design。

\lecturefigure{slide-34.jpg}{Chain-of-thought discussion：prompting 的持久性与开放问题}{官方课件第 34 页}

\begin{knowledgebox}{读图：把已知结论与未知机制分开}
课堂证据支持 CoT 在多个 benchmark、语言和模型上有效，并显示规模门槛；尚未解决的是 pretraining 中哪些经验使模型学会使用 reasoning traces、这些 traces 与内部计算有何关系、哪些任务会从 few-shot CoT 泛化。Discussion 页的价值正是把成功结果后面的 unknowns 留在视野中。
\end{knowledgebox}

\teachervoice{有人问 CoT 是否只是额外 token 带来的 compute。讲者提到 equal-length filler 控制无效，并判断语言在引导 reasoning 中很重要；有人问先给答案再解释，他说通常更差。这些是重要 negative controls，但他仍把 ``why CoT works'' 明确称为开放问题。}

\subsection{整堂课的结论}

课件将前后两部分压成三层：能力可以随规模在任务 metric 上呈现阈值；prompting、fine-tuning 和 data quality 会移动可测边界；CoT 与 self-consistency 说明 inference-time method 也必须与 base-model capability 匹配。

\lecturefigure{slide-35.jpg}{课程结论：Scaling、CoT 与 self-consistency 的统一视角}{官方课件第 35 页}

\begin{importantbox}{整堂课最值得带走的六点}
\begin{enumerate}
  \item Smooth loss scaling 与 thresholded task performance 可以同时成立。
  \item Emergence 的判断依赖 model family、prompt、metric 与 baseline。
  \item Better data、architecture 与 post-training 能把经验阈值向小模型移动。
  \item CoT 通过中间 token 提供 decomposition 与 sequential computation，但不保证解释忠实。
  \item BBH 与 multilingual CoT 显示收益跨任务、跨语言，却仍具有明显异质性。
  \item Self-consistency 把 inference samples 变成 accuracy，必须显式报告 cost 与 error correlation。
\end{enumerate}
\end{importantbox}

\subsection{讲者给出的研究方向}

2023 年 1 月的现实是，纯粹扩大 foundation model 主要由少数工业实验室承担。讲者因此把学生可进入的方向放在 better prompting、capability characterization、applications、benchmarks 和 compute-efficient methods 上：这些工作既能更好利用现有模型，也能检验 scaling narrative 的边界。

\lecturefigure{slide-36.jpg}{Looking forward：讲者的个人研究兴趣}{官方课件第 36 页}

\begin{knowledgebox}{读图：五个方向分别解决哪一层问题}
Scaling 研究扩大 base capability；better prompting 改善 elicitation；characterization 建立能力地图；applied work 把模型放入真实 workflow；benchmarks 提供不会迅速饱和的测量；compute-efficient methods 则降低训练与推理门槛。它们不是相互替代，而是从能力、接口、测量、应用和资源五层共同推进。
\end{knowledgebox}

\teachervoice{讲者收尾时说，现有 benchmark 被解决得很快，需要更难的新评测；同时应研究更 compute-efficient 的方法，让并非只有工业界才能使用更强模型。这一结尾把 scaling 的收益与访问成本放在同一张研究议程中。}

\subsection{本章小结}

Prompt engineering 的形式可能变化，但明确目标、分解任务与验证结果不会自动消失。课程的成熟之处在于：行为结果很强，机制解释仍谨慎；扩展模型值得研究，数据、评测和效率同样决定能力能否被发现、复现与部署。

